% =====================================================================
%  Capítulo 3 — Arquitetura do sistema
% =====================================================================
\chapter{Arquitetura do sistema}
\label{cap:arquitetura}
\rubricaJustificativa

\section{Visão em camadas}

A Figura~\ref{fig:arquitetura} situa o sistema da AP2 nas quatro camadas
de uma arquitetura IoT \cite{alfuqaha2015}. A camada de rede, que na AP1
era o cabo USB, agora é um enlace BLE entre o nó e um central.

\begin{itemize}
  \item \textbf{Percepção:} DHT11 (temperatura e umidade do ar) e BH1750
    (luz em lux, por I\textsuperscript{2}C). Os LEDs de aviso e o módulo
    fotorresistor saíram.
  \item \textbf{Processamento na borda:} o ESP32 mede, tenta de novo
    quando um sensor falha, integra a luz no tempo e monta o pacote de
    \SI{19}{\byte}. Ele não decide mais nada sobre a planta: a comparação
    com a espécie passou para o aplicativo (\DP{3}).
  \item \textbf{Rede:} BLE. O nó é periférico GATT; o central assina a
    notificação, recebe um pacote a cada \SI{300}{\second} e confirma
    cada um. Não há roteador, \emph{broker} nem internet entre o nó e
    quem guarda o dado.
  \item \textbf{Aplicação:} na campanha, o gravador em Python, que grava
    o CSV, e o \texttt{qualidade.py}, que o analisa. No produto, o
    aplicativo FitoSfera no celular, seguindo o mesmo contrato.
\end{itemize}

\begin{figure}[H]
\centering
\begin{tikzpicture}[
  font=\sffamily\scriptsize,
  caixa/.style={draw, rounded corners=2pt, align=center, minimum height=0.9cm, inner sep=3pt, fill=white},
  futuro/.style={caixa, dashed},
  camada/.style={font=\sffamily\scriptsize\bfseries, anchor=west, align=left},
  seta/.style={-{Stealth[length=2mm]}, thick},
]
  \foreach \y/\nome in {0/Percepção, 1.6/Processamento\\(borda), 3.2/Rede, 4.8/Aplicação}{
    \fill[gray!8] (-0.2,\y-0.65) rectangle (13.6,\y+0.65);
    \node[camada] at (-0.1,\y) {\nome};
  }
  \node[caixa, text width=2.4cm] (dht) at (3.6,0) {DHT11\\T e UR, GPIO18};
  \node[caixa, text width=2.6cm] (bh)  at (6.8,0) {BH1750 (0x23)\\luz, I\textsuperscript{2}C 21/22};
  \node[caixa, text width=4.6cm] (esp) at (5.2,1.6) {ESP32 DevKit V1 a \SI{80}{\mega\hertz}\\mede, integra a luz, monta 19 bytes};
  \node[caixa, text width=4.6cm] (ble) at (5.2,3.2) {BLE: notificação a cada \SI{300}{\second}\\confirmação \texttt{[0x02,} \texttt{seq]}};
  \node[caixa, text width=3.0cm] (grav) at (3.2,4.8) {gravador.py (PC)\\\texttt{mensagens.csv}};
  \node[caixa, text width=2.2cm] (qual) at (7.0,4.8) {qualidade.py};
  \node[futuro, text width=2.8cm] (app) at (11.6,4.8) {aplicativo FitoSfera\\no celular (AP3)};
  \node[futuro, text width=2.8cm] (ser) at (11.6,1.6) {serial USB\\CSV de depuração};
  \draw[seta] (dht) -- (esp);
  \draw[seta] (bh) -- (esp);
  \draw[seta] (esp) -- (ble);
  \draw[seta] (ble) -- (grav);
  \draw[seta] (grav) -- (qual);
  \draw[seta, dashed] (ble.east) -| (app.south);
  \draw[seta, dashed] (esp) -- (ser);
\end{tikzpicture}
\caption{Arquitetura do FitoSfera na AP2: o que rodou na campanha (contínuo) e o que fica para depois ou só serve à depuração (tracejado)}
\label{fig:arquitetura}
\fonte{Elaborada pela equipe.}
\end{figure}

\section{Componentes e decisões de projeto}
\label{sec:componentes}

As caixas \DP{1} a \DP{5} vêm da AP1 e foram atualizadas para o sistema
de agora. As caixas \DP{6} a \DP{10} são do \moduloDoCiclo: o que mudou
na arquitetura por causa do enlace. A escolha da tecnologia em si está
na \DP{11}, no Capítulo~\ref{cap:conectividade}.

\begin{decisao}{Sensor de temperatura e umidade do ar (mantida)}
\alternativas{DHT11, DHT22 e BME280 (I\textsuperscript{2}C).}
\criterio{Estar disponível para a equipe, custo e resolução suficiente
  para perceber tendências num ambiente que muda em escala de horas.}
\escolha{DHT11, como na AP1.}
\evidencia{Resolução de \SI{1}{\celsius} e 1\,\% UR e incerteza de
  $\pm$\SI{2}{\celsius} \cite{aosongdht11} (\folhadedados). Com uma
  leitura a cada \SI{300}{\second}, o intervalo mínimo de \SI{2}{\second}
  da biblioteca, que na AP1 coincidia com o período (falha 2 da
  Seção~5.3 da AP1), deixou de ser um risco.}
\consequencia{A leitura continua ocupando cerca de \SI{25}{\milli\second}
  (Apêndice~\ref{ap:barramento}). O BME280 é a troca prevista: entra no
  mesmo barramento I\textsuperscript{2}C do BH1750, sem pino novo.}
\end{decisao}

\begin{decisao}{Canal de luz: BH1750 em I\textsuperscript{2}C}
\alternativas{(a) manter o fotorresistor numa entrada digital, como na
  AP1; (b) ler o fotorresistor pelo ADC1 (GPIO32 ou GPIO33);
  (c) sensor de luz calibrado em I\textsuperscript{2}C (BH1750).}
\criterio{O aplicativo precisa saber \emph{quanta} luz a planta recebeu
  num dia, e não se o ambiente está claro ou escuro; o limiar tem de ser
  de software.}
\escolha{(c), BH1750 (módulo GY-302) no endereço \texttt{0x23}, no modo
  de uma medida, com MTreg = 31.}
\evidencia{A AP1 mostrou que (a) lê uma tensão contínua numa entrada
  digital, numa faixa de \SI{1.65}{\volt} sem garantia da folha de dados.
  (b) daria um número, mas não em lux: o fotorresistor não tem curva
  conhecida e o ADC do ESP32 é pouco linear nos extremos. O BH1750 entrega
  lux diretamente, com incerteza de $\pm$20\,\% \cite{rohmbh1750}
  (\folhadedados). Na sensibilidade de fábrica (MTreg = 69) ele satura
  perto de \SI{54612}{lux}; com MTreg = 31 o teto sobe para cerca de
  \SI{121557}{lux} ($54612 \times 69/31$), o que cobre sol pleno.}
\consequencia{Cada medida espera cerca de \SI{100}{\milli\second} de
  conversão com \texttt{delay()}, aceito porque acontece uma vez a cada
  \SI{300}{\second}. A resolução em luz fraca cai na mesma proporção
  ($69/31 \approx 2{,}2$ vezes).}
\end{decisao}

\begin{decisao}{Onde fica a decisão sobre a planta}
\alternativas{(a) comparador no hardware; (b) limiares no
  \emph{firmware}, com LEDs no nó (a escolha da AP1); (c) nó que só mede,
  com a decisão no aplicativo.}
\criterio{O limiar tem de ser o da espécie cadastrada, diferente para
  cada planta, e tem de poder mudar sem regravar o nó.}
\escolha{(c). Na AP2 os três LEDs e as constantes \texttt{TEMP\_ALTA} e
  \texttt{UMID\_BAIXA} saíram do \emph{firmware}.}
\evidencia{Na AP1 a equipe já tinha declarado que os LEDs eram de
  demonstração e que o limiar fixo na compilação valia para o jardim
  inteiro. Com o BLE, o aplicativo recebe as leituras e conhece a
  espécie; mandar os limiares de volta para o nó só serviria para acender
  um LED que ninguém olha.}
\consequencia{Sem central, não há aviso nenhum. É o preço de (c), e é
  aceitável porque a planta muda em horas e o celular volta a conectar
  quando o usuário chega perto do jardim. Na AP3, com a IA embarcada,
  a decisão pode voltar em parte para o nó (Seção~\ref{sec:proximo-ciclo}).}
\end{decisao}

\begin{decisao}{Temporização sem bloqueio (atualizada)}
\alternativas{\texttt{delay(300000)} no fim do \texttt{loop()}, ou
  comparação com \texttt{millis()}.}
\criterio{O laço tem de atender, a qualquer momento, a assinatura da
  notificação, o comando de medir agora e a confirmação do central.}
\escolha{Comparação com \texttt{millis()} contra um instante
  \texttt{proximaLeitura} (Listagem~\ref{lst:laco}). Um pedido do central
  fura a fila e mede na hora.}
\evidencia{Com \texttt{delay()}, um pedido do central esperaria até
  \SI{300}{\second}; com a comparação, a resposta sai no próximo giro do
  laço, e \RF{08} pede no máximo \SI{2}{\second}. Os retornos do BLE
  chegam por \emph{callbacks} do NimBLE, que só ligam uma bandeira; o
  trabalho é feito no \texttt{loop()}.}
\consequencia{A única espera bloqueante que sobrou são os
  \SI{100}{\milli\second} da conversão do BH1750 (\DP{2}), uma vez por
  ciclo.}
\end{decisao}

\begin{decisao}{Escolha dos pinos (atualizada)}
\alternativas{Qualquer GPIO de uso geral da placa.}
\criterio{Evitar os pinos que a folha de dados proíbe ou desaconselha
  \cite{espressif2023}: GPIO6 a 11 (\emph{flash}), GPIO0, 2, 5, 12 e 15
  (inicialização), GPIO1 e 3 (UART0) e GPIO34 a 39 (só entrada).}
\escolha{GPIO18 (DHT11) e GPIO21/22 (SDA/SCL do BH1750).}
\evidencia{O conflito previsto na AP1 era o I\textsuperscript{2}C padrão
  cair em cima dos LEDs vermelho e amarelo. Com a saída dos LEDs, o par
  21/22 ficou livre e o barramento usa os pinos padrão, sem
  \texttt{Wire.begin()} em pinos alternativos (\CT{10}).}
\consequencia{Três pinos usados, nenhum restrito. Qualquer sensor
  I\textsuperscript{2}C novo entra no mesmo par, desde que o endereço seja
  diferente de \texttt{0x23}.}
\end{decisao}

\begin{decisao}{O que conta como entregue: confirmação da aplicação}
\alternativas{(a) considerar entregue quando \texttt{notify()} volta;
  (b) usar indicação em vez de notificação, que a pilha BLE confirma;
  (c) notificação mais uma confirmação escrita pela aplicação,
  \texttt{[0x02,} \texttt{seq]}.}
\criterio{\enquote{Entregue} tem de significar \enquote{gravado por quem
  guarda o dado}, e não \enquote{saiu pelo rádio}.}
\escolha{(c). O nó guarda a sequência do pacote em espera e só zera o
  contador de segundos sem sucesso quando chega a confirmação dessa
  sequência. O gravador grava a linha com \texttt{fsync} e só depois
  confirma (Listagem~\ref{lst:gravador}).}
\evidencia{Uma notificação não tem confirmação no protocolo: a volta de
  \texttt{notify()} só prova que o pacote foi para a fila do rádio. Uma
  indicação é confirmada pela pilha do central, não pelo aplicativo, que
  pode cair depois de receber e antes de gravar. No ensaio com o nRF
  Connect, sem a escrita de \texttt{02} \texttt{NN}, o contador continuou
  crescendo de pacote em pacote (\CT{14}).}
\consequencia{O campo \texttt{seg\_sem\_sucesso} de cada pacote diz há
  quanto tempo o nó não tem prova de entrega. Perto de \SI{300}{\second},
  tudo certo; muito acima, a confirmação anterior se perdeu ou houve
  lacuna. O \texttt{qualidade.py} conta esses casos.}
\end{decisao}

\begin{decisao}{Formato do pacote: binário de 19 bytes}
\alternativas{(a) JSON em texto, como o contrato de dados da disciplina;
  (b) estrutura binária fixa, \emph{little-endian}.}
\criterio{Caber numa única notificação sem depender de negociar o MTU,
  e ser decodificável sem ambiguidade por um aplicativo em TypeScript.}
\escolha{(b), 19 bytes (Tabela~\ref{tab:contrato-dados}).}
\evidencia{Com o MTU padrão do ATT, de \SI{23}{\byte}, cabem
  \SI{20}{\byte} de dados numa notificação. O JSON equivalente, com as
  mesmas oito grandezas e sem espaços, tem \SI{122}{\byte}: exigiria
  negociar um MTU maior, o que nem todo celular faz, ou fragmentar a
  mensagem. Os três pacotes de exemplo do README foram decodificados pelo
  decodificador em TypeScript e pelo \texttt{protocolo.py} sem divergência
  (\CT{12}).}
\consequencia{O pacote não se descreve sozinho: quem recebe precisa do
  contrato. Uma mudança de campo exige versão nova do contrato (hoje,
  versão 2) e atualização do decodificador nos dois lados.}
\end{decisao}

\begin{decisao}{O que o nó guarda enquanto ninguém escuta}
\alternativas{(a) nada; (b) uma fila circular de leituras na RAM,
  entregue quando o central volta; (c) só a luz, como integral
  acumulada desde que o nó ligou.}
\criterio{A grandeza que o aplicativo soma ao longo do dia (luz) não
  pode ter buraco; as que ele só compara com um limiar (temperatura e
  umidade) toleram.}
\escolha{(c). A cada medida, o nó soma a área do trapézio entre a leitura
  atual e a anterior; intervalos de mais de \SI{15}{\minute} sem leitura
  de luz não entram na soma.}
\evidencia{Num ensaio com o nRF Connect, depois de cerca de
  \SI{55}{\minute} desconectado, o primeiro pacote trouxe de uma vez a
  luz do período inteiro (\CT{20}). Como o aplicativo soma diferenças
  entre acumulados, um pacote repetido não conta luz em dobro e um
  perdido não deixa buraco. Uma fila de um dia (288 leituras $\times$
  \SI{19}{\byte} $\approx$ \SI{5.5}{\kilo\byte}) caberia na RAM, mas
  exigiria um protocolo de entrega em lote que a equipe não implementou
  neste ciclo.}
\consequencia{Temperatura e umidade medidas sem central se perdem. A
  sequência continua contando, então o conjunto de dados mostra quantas
  leituras faltaram, mas não os valores. É por isso que, na campanha, o
  gravador precisa ficar ligado o tempo todo (Seção~\ref{sec:caminho}).}
\end{decisao}

\begin{decisao}{Período de amostragem: \SI{300}{\second}}
\alternativas{\SI{2}{\second} (AP1), \SI{60}{\second} e \SI{300}{\second}.}
\criterio{Ser bem mais rápido que a resposta da planta, que é de horas,
  sem gastar rádio e sem inflar o conjunto de dados.}
\escolha{\SI{300}{\second} (\texttt{INTERVALO\_MS}), com leitura extra a
  pedido do central.}
\evidencia{São 288 leituras e $288 \times \SI{19}{\byte} =
  \SI{5472}{\byte}$ por dia, doze amostras por hora. A \SI{2}{\second}
  seriam 43.200 notificações por dia, sem informação nova sobre a planta.}
\consequencia{A luz é amostrada a cada \SI{5}{\minute}, o que torna a
  integral grosseira em dias de nuvens passando. Um evento mais curto que
  \SI{5}{\minute} (uma sombra passageira) pode não aparecer no conjunto
  de dados.}
\end{decisao}

\begin{decisao}{Quem dá a hora da leitura}
\alternativas{(a) NTP no nó, o que exige WiFi; (b) relógio de tempo real
  externo; (c) o central carimba a hora na chegada.}
\criterio{Toda leitura do conjunto de dados precisa de hora real, para
  que a série tenha ciclo diário e para rotular eventos.}
\escolha{(c). O gravador grava \texttt{ts\_recebido} com o relógio do
  computador, sincronizado pelo sistema operacional; o aplicativo fará o
  mesmo com o relógio do celular.}
\evidencia{Com BLE não há rede no nó, então (a) não se aplica; (b) custaria
  um módulo e uma bateria para corrigir um erro que, com notificação, é do
  tamanho do intervalo de conexão (dezenas de milissegundos). O nó mede e
  notifica no mesmo giro do laço.}
\consequencia{Nenhuma leitura sai com hora do nó: o equivalente a
  \enquote{\texttt{ts} \texttt{null}} da disciplina vale para 100\,\% dos
  pacotes, e é tratado como decisão, não como falha (\CT{19}). Uma leitura
  feita por \texttt{READ}, sem notificação, pode ter até \SI{300}{\second}
  de idade; o gravador não usa \texttt{READ}.}
\end{decisao}

\section{Mapa de barramento}
\label{sec:barramento}

O nó tem dois sensores e o enlace. O DHT11 continua no barramento de um
fio próprio, no GPIO18, e é o dispositivo mais lento, com cerca de
\SI{25}{\milli\second} por leitura. O BH1750 está sozinho no
I\textsuperscript{2}C, no endereço \texttt{0x23} (pino ADDR solto), nos
pinos padrão 21/22, que a saída dos LEDs liberou (\DP{5}). Na
inicialização o \emph{firmware} varre o barramento e lista os endereços
na serial; se o \texttt{0x23} não responde, o nó avisa e segue sem luz.
O mapa completo está no Apêndice~\ref{ap:barramento}.

\section{Orçamento de energia}
\label{sec:energia}

O nó continua sempre ligado, num carregador USB de \SI{5}{\volt}, agora
com o BLE anunciando ou conectado o tempo todo e a CPU a
\SI{80}{\mega\hertz}, o mínimo com BLE. Sem os LEDs, o consumo estimado
cai de \SI{66}{\milli\ampere} (AP1) para cerca de
\SI{46}{\milli\ampere} em média, ou cerca de
\SI{1114}{\milli\ampere\hour} por dia. Quem domina é o ESP32 acordado com
o rádio, que responde por cerca de 97\,\% da conta; os sensores ficam
abaixo de 2\,\%. A linha do rádio e a corrente total estão detalhadas no
Apêndice~\ref{ap:energia} e na Seção~\ref{sec:custo-radio}. A equipe não
tem medidor USB, então esses números continuam sendo estimativa.

\section{Lista de materiais}
\label{sec:bom}

\begin{table}[H]
\centering
\caption{Lista de materiais do protótipo}
\label{tab:bom}
\footnotesize
\begin{tabular}{@{}L{6.0cm}C{1.2cm}R{2.2cm}R{2.2cm}L{2.2cm}@{}}
\toprule
\textbf{Item} & \textbf{Qtd.} & \textbf{Unitário (R\$)} & \textbf{Total (R\$)} & \textbf{Origem} \\
\midrule
Kit eletrônica iniciante ESP32 & 1 & 180,00 & 180,00 & Mercado Livre \\
\multicolumn{5}{@{}l}{\emph{Itens do kit usados no protótipo:}} \\
\quad Placa ESP32 DevKit V1 (ESP32-WROOM-32) & 1 & --- & --- & kit \\
\quad Módulo DHT11 (3 pinos, \emph{pull-up} na placa) & 1 & --- & --- & kit \\
\quad Protoboard e fios macho-macho (9 em uso) & 1 & --- & --- & kit \\
\quad Cabo USB & 1 & --- & --- & kit \\
Módulo BH1750 (GY-302) & 1 & 21,04 & 21,04 & Mercado Livre \\
Carregador USB \SI{5}{\volt}, \SI{1}{\ampere} ou mais & 1 & --- & --- & da equipe \\
\midrule
\multicolumn{3}{@{}r}{\textbf{Total}} & \textbf{201,04} & \\
\bottomrule
\end{tabular}
\fonte{Elaborada pela equipe.}
\end{table}

Saíram do protótipo da AP1 os três LEDs, os três resistores de
\SI{220}{\ohm} e o módulo fotorresistor. O gravador roda num notebook Acer Aspire A515-57
que a equipe já tinha (Linux Mint 22.3, Python 3.12.3), com Bluetooth.
